% Modelado supervisado del proyecto PSIM (laboratorio de regresion)
% Compilar desde la raiz:
% latexmk -pdf -interaction=nonstopmode -halt-on-error -outdir=build/reports reports/informe.tex

\documentclass[11pt,a4paper]{article}

\usepackage[utf8]{inputenc}
\usepackage[T1]{fontenc}
\usepackage[spanish]{babel}
\usepackage{lmodern}
\usepackage{geometry}
\geometry{left=2.6cm,right=2.6cm,top=2.8cm,bottom=2.6cm,headheight=1.4cm}
\usepackage{setspace}
\onehalfspacing

\usepackage{hyperref}
\hypersetup{colorlinks=true,linkcolor=blue,citecolor=blue,urlcolor=blue}
\usepackage{graphicx}
\usepackage{fancyhdr}
\usepackage{amsmath,amssymb}
\usepackage{booktabs}
\usepackage{caption}
\usepackage{float}
\usepackage{longtable}
\usepackage{array}

\pagestyle{fancy}
\fancyhf{}
\fancyhead[L]{\includegraphics[height=0.95cm]{Escuela_Rosario_logo.png}}
\fancyfoot[C]{\thepage}
\renewcommand{\headrulewidth}{0pt}

\makeatletter
\providecommand{\input@path}{}
\edef\input@path{{reports/}{reports/tablas/}{./}}
\makeatother
\graphicspath{{reports/}{figures/modelado_proyecto/}{results/modelado_didactico/}{./}}

\input{lab005_cifras}

\newcommand{\tituloTrabajo}{Modelado supervisado con regresion multiple}
\newcommand{\subtituloTrabajo}{Estimacion de la fraccion pulmonar de referencia a partir de descriptores de la segmentacion automatica}

\begin{document}

\thispagestyle{empty}
\begin{titlepage}
    \centering
    {\LARGE \textbf{\tituloTrabajo}}\\[0.6cm]
    {\large \subtituloTrabajo}\\[2cm]
    {\large \textbf{Juanita Trujillo Narvaez}. Cod: 1000099467}\\
    {\large \textbf{Karol Isabella Martinez Villarreal}. Cod: 1000099682}\\[1cm]
    {\large Ingenieria Biomedica}\\[0.4cm]
    {\large PSIM - 2026II - 80}\\
    Docente: Ing. Pablo Eduardo Caicedo-Rodriguez. Ph.D.
    \vfill
    Fecha de entrega: \today
\end{titlepage}

\tableofcontents
\newpage

\begin{abstract}
Este laboratorio extiende la tabla de caracteristicas del analisis exploratorio
con un modelo supervisado. Se pregunta si los descriptores de la segmentacion
automatica sin realce y con umbral adaptativo permiten estimar la fraccion
pulmonar que mide la mascara de referencia trazada por personas, en pacientes
que el modelo no uso para ajustar sus parametros. Se comparan la regresion
multiple de grado 1 y su extension polinomica de grado 2 con tres predictores,
frente a la media de entrenamiento como referencia sencilla. El modelo de
\ModElegidoNombre{}, elegido por RMSE de validacion, alcanzo en la prueba de
\ModNTest{} pacientes un RMSE de \ModRMSETest{} frente a \ModRMSETestBase{} de la
referencia ($R^2 = \ModRTest{}$). El error fue del mismo orden en los dos
origenes y sin sesgo medio apreciable, aunque Montgomery aporta pocos casos.
Se incluye ademas la practica comun con datos sinteticos, en la que se
desarrollan las rutas continua y binaria.
\end{abstract}

\newpage

\section{Pregunta y continuidad con el analisis exploratorio}

El analisis exploratorio comparo nueve estrategias de segmentacion clasica y
eligio \texttt{base\_adap} (sin realce, umbral adaptativo) con un criterio
fijado de antemano. Para cada imagen con mascara de referencia, la tabla de
caracteristicas contiene cinco descriptores morfologicos calculados sobre la
mascara obtenida y los mismos descriptores sobre la mascara manual.

Ese analisis describe cuanto se parecen ambos valores. Este laboratorio
formula la pregunta predictiva:

\begin{quote}
\emph{?`Los descriptores de la segmentacion automatica permiten estimar la
fraccion pulmonar de la mascara de referencia en pacientes que el modelo no ha
utilizado para ajustar sus parametros?}
\end{quote}

\begin{table}[H]
\centering
\small
\caption{Contrato de modelado}
\begin{tabular}{@{}ll@{}}
\toprule
Unidad independiente & paciente (\texttt{grupo\_id}) \\
Salida & \texttt{ref\_fraccion\_area}, fraccion adimensional del campo irradiado \\
Referencia & mascara pulmonar trazada por personas (NLM) \\
Predictores & \texttt{fraccion\_area}, \texttt{energia\_wavelet\_d1}, \texttt{simetria\_izq\_der} \\
Estrategia & \texttt{base\_adap}, imagenes elegibles \\
Metrica primaria & RMSE, fijada antes de comparar \\
Ruta binaria & no aplica \\
\bottomrule
\end{tabular}
\end{table}

La ruta binaria no aplica al proyecto. El proyecto no formula la deteccion de
tuberculosis, y su contrato de entrada declara la etiqueta clinica como
variable de control. No se crea una segunda pregunta para cumplir la
actividad; ambas rutas se desarrollan en la practica comun.

\section{Datos y particiones}

Se conserva la particion congelada del laboratorio anterior (70/15/15 por
unidad independiente, estratificada por origen y clase). Los descriptores de
validacion y prueba se extrajeron con el mismo codigo del analisis
exploratorio, restringido a \texttt{base\_adap} y a las imagenes con mascara de
referencia. Las filas de prueba se generaron \textbf{despues} de congelar el
protocolo.

Una imagen es elegible si supero el control de calidad y su segmentacion
produjo los dos campos pulmonares. La regla se aplica igual en las tres
particiones. Las imagenes de un mismo paciente se promedian en una sola fila.

\begin{table}[H]
\centering
\small
\caption{Conteos por particion}
\begin{tabular}{@{}lrrr@{}}
\toprule
\textbf{Particion} & \textbf{Con mascara} & \textbf{Elegibles} & \textbf{Pacientes} \\
\midrule
Entrenamiento & \ModMascTrain & \ModElegTrain & \ModNTrain \\
Validacion & \ModMascVal & \ModElegVal & \ModNVal \\
Prueba & \ModMascTest & \ModElegTest & \ModNTest \\
\bottomrule
\end{tabular}
\end{table}

Los predictores se seleccionaron solo con entrenamiento, segun su asociacion
con el objetivo, la consistencia de signo entre origenes y la redundancia entre
ellos. La asociacion de Spearman con el objetivo fue \ModRhoFA{} para
\texttt{fraccion\_area}, \ModRhoEW{} para \texttt{energia\_wavelet\_d1} y
\ModRhoSI{} para \texttt{simetria\_izq\_der}. La Tabla~\ref{tab:decisiones}
resume las decisiones.

\begingroup
\renewcommand{\arraystretch}{1.15}
\phantomsection\label{tab:decisiones}
\input{lab005_decisiones}
\endgroup

\section{Metodo}

Sea $z_j = (x_j - \bar{x}_j)/s_j$ el predictor $j$ estandarizado con la media y
la desviacion estandar de entrenamiento. Los candidatos son:

\begin{align}
\text{Grado 1:}\quad \hat{y} &= \beta_0 + \sum_{j=1}^{3} \beta_j z_j, \\
\text{Grado 2:}\quad \hat{y} &= \beta_0 + \sum_{j=1}^{3} \beta_j z_j
  + \sum_{j=1}^{3}\sum_{k=j}^{3} \beta_{jk}\, z_j z_k .
\end{align}

El modelo de grado 2 tiene nueve terminos mas el intercepto. Es no lineal en los
predictores y lineal en los coeficientes, que se estiman por minimos
cuadrados. La imputacion por mediana, el escalamiento y la expansion se ajustan
dentro de un \emph{pipeline} solo con entrenamiento.

La referencia sencilla es la media de entrenamiento; la mediana se reporta como
referencia adicional. La metrica primaria es
$\mathrm{RMSE} = \sqrt{\tfrac{1}{n}\sum_i (y_i - \hat{y}_i)^2}$, en las mismas
unidades de la salida; se reportan tambien MAE y $R^2$. El candidato se elige
por menor RMSE de validacion. Despues, el protocolo se congela y la prueba se
evalua una sola vez, reconstruyendo el ajuste desde entrenamiento y verificando
que los coeficientes coinciden con los congelados.

\section{Practica comun con datos sinteticos}

La practica usa 240 sujetos simulados con tres descriptores y dos salidas
generadas por mecanismos conocidos. No son imagenes reales ni diagnosticos.

En la ruta continua, el modelo de \DidModeloReg{} reconoce el termino
cuadratico del mecanismo. Obtuvo un RMSE de \DidRMSEValReg{} en validacion y
\DidRMSETestReg{} en prueba, frente a \DidRMSETestBase{} de la referencia. El
RMSE no llega a cero por el ruido aleatorio incluido en la simulacion.

En la ruta binaria, \DidModeloBin{} alcanzo en prueba un ROC-AUC de
\DidAUCTest{} y un AP de \DidAPTest{}, frente a una prevalencia de
\DidPrevTest{}. Con el umbral de referencia 0.5, la sensibilidad fue
\DidSensTest{} y la especificidad \DidEspTest{}. Una probabilidad correcta no
clasifica todos los casos, porque la clase se genera por un sorteo de
Bernoulli. Cambiar el umbral o $C$ despues de mirar la prueba la convertiria en
otra validacion.

\input{lab005_didactico_reg}
\input{lab005_didactico_bin}

\section{Resultados del proyecto}

\subsection{Seleccion en validacion}

El modelo de grado 2 obtuvo el menor RMSE de validacion (\ModRMSEValDos{}),
frente a \ModRMSEValUno{} del grado 1 y \ModRMSEValBase{} de la media. La
ventaja de la extension polinomica es modesta: la mayor parte del poder
predictivo ya esta en el modelo lineal. La matriz de diseno tiene rango
completo (\ModRangoDiseno{}) y numero de condicion \ModCondicion{}, de modo que
los coeficientes son identificables.

\input{lab005_validacion}

\begin{figure}[H]
\centering
\includegraphics[width=\textwidth]{validacion_prediccion_residuos.png}
\caption{Validacion, \ModNVal{} pacientes: referencia frente a prediccion y
residuos frente a prediccion, por origen. Unidades: fraccion adimensional.}
\end{figure}

\begin{figure}[H]
\centering
\includegraphics[width=\textwidth]{validacion_residuos_predictores.png}
\caption{Validacion: residuos frente a cada predictor.}
\end{figure}

\input{lab005_coeficientes}

Los coeficientes se expresan por desviacion estandar de entrenamiento.
\texttt{fraccion\_area} domina, y los terminos cuadraticos y las interacciones
son pequenos. Son asociaciones observacionales: no indican que un descriptor
cause la fraccion pulmonar.

\subsection{Evaluacion unica en prueba}

\input{lab005_prueba}

En la prueba de \ModNTest{} pacientes, el modelo alcanzo un RMSE de
\ModRMSETest{} (MAE \ModMAETest{}), frente a \ModRMSETestBase{} de la media de
entrenamiento: una reduccion del \ModMejoraTest\,\% ($R^2 = \ModRTest{}$). El
desempeno coincide con el de validacion ($R^2 = \ModRVal{}$), sin optimismo
visible por la seleccion. El RMSE fue \ModRMSETestSH{} en Shenzhen
($n=\ModNTestSH{}$) y \ModRMSETestMC{} en Montgomery ($n=\ModNTestMC{}$). El
residuo medio fue $\ModResTestSH{}$ y $\ModResTestMC{}$, respectivamente.

\begin{figure}[H]
\centering
\includegraphics[width=\textwidth]{prueba_prediccion_residuos.png}
\caption{Prueba, \ModNTest{} pacientes: referencia frente a prediccion y residuos.}
\end{figure}

\section{Discusion y limites}

\begin{itemize}
    \item \textbf{Regresion a la media.} Las fracciones extremas se sobreestiman
    (las bajas) o se subestiman (las altas). El modelo es menos confiable en
    esos casos.
    \item \textbf{Efecto de centro.} La fraccion de referencia media en
    entrenamiento fue \ModRefMediaMC{} en Montgomery y \ModRefMediaSH{} en
    Shenzhen. La referencia constante falla sistematicamente en Montgomery; el
    modelo corrige ese desfase con los descriptores. Montgomery aporta solo
    \ModNTestMC{} pacientes en prueba, de modo que la comparacion entre centros
    es imprecisa. No se ha verificado el transporte a un tercer centro.
    \item \textbf{Denominador compartido.} \texttt{fraccion\_area} y el objetivo
    se normalizan por el mismo campo irradiado. La asociacion no es circular,
    porque provienen de mascaras distintas, pero comparten ese componente.
    \item \textbf{Segmentaciones fallidas.} El modelo no aplica cuando la
    segmentacion automatica produce una sola componente.
    \item \textbf{Soporte de los predictores.} \texttt{simetria\_izq\_der} es muy
    asimetrica; en sus valores extremos hay pocos casos y el termino cuadratico
    extrapola.
    \item \textbf{Incertidumbre.} Se reportan estimaciones puntuales; no se
    calcularon intervalos por bootstrap.
\end{itemize}

\section{Reproduccion y comprobaciones}

Todo el flujo se reconstruye desde los datos originales con:

\begin{verbatim}
uv sync --locked
uv run ruff check .
uv run ruff format --check .
uv run pytest
uv run python scripts/reproduce.py
latexmk -pdf -interaction=nonstopmode -halt-on-error
        -outdir=build/reports reports/informe.tex
\end{verbatim}

El entorno usa Python 3.13 con dependencias bloqueadas en \texttt{uv.lock}. El
contrato registra el hash SHA-256 de la tabla de entrenamiento, y el protocolo
congelado el de validacion; la prueba falla si alguno cambia. Una segunda
ejecucion reprodujo exactamente las metricas de prueba.

Las pruebas automaticas verifican que:
\begin{itemize}
    \item un paciente que cruce particiones produce error;
    \item el objetivo, los identificadores y las columnas derivadas de la
    mascara manual no pueden ser predictores;
    \item la imputacion aprende solo de entrenamiento;
    \item un objetivo faltante se excluye y no se imputa;
    \item las metricas coinciden con valores calculados a mano;
    \item dos ajustes iguales producen los mismos coeficientes.
\end{itemize}

\section{Conclusion}

En la prueba de \ModNTest{} pacientes, la regresion polinomica de grado
\ModGrado{} sobre tres descriptores de la segmentacion automatica estimo la
fraccion pulmonar de referencia con un RMSE de \ModRMSETest{}, frente a
\ModRMSETestBase{} de la media de entrenamiento. El error fue del mismo orden en
ambos origenes, sin sesgo medio apreciable. El modelo regresa a la media en
fracciones extremas, no se ha verificado su transporte a otro centro y no
aplica cuando la segmentacion falla.

\end{document}